\documentclass{article}

\usepackage{fancyhdr}
\usepackage{extramarks}
\usepackage{amsmath}
\usepackage{amsthm}
\usepackage{amsfonts}
\usepackage{tikz}
\usepackage[plain]{algorithm}
\usepackage{algpseudocode}
\usepackage{listings}
\usepackage{xcolor}
\usepackage{float}

\usetikzlibrary{automata,positioning}

\newcommand{\overbar}[1]{\mkern 1.5mu\overline{\mkern-1.5mu#1\mkern-1.5mu}\mkern 1.5mu}

%
% Basic Document Settings
%

\topmargin=-0.45in
\evensidemargin=0in
\oddsidemargin=0in
\textwidth=6.5in
\textheight=9.0in
\headsep=0.25in

\linespread{1.1}

\pagestyle{fancy}
\lhead{Yousef Alaa Awad}
\chead{\hmwkClass\: \hmwkTitle}
\rhead{\firstxmark}
\lfoot{\lastxmark}
\cfoot{\thepage}

\renewcommand\headrulewidth{0.4pt}
\renewcommand\footrulewidth{0.4pt}

\setlength\parindent{0pt}

%
% Create Problem Sections
%

\setcounter{secnumdepth}{0}

\newcommand{\hmwkTitle}{Assignment\ \#5}
\newcommand{\hmwkClass}{CDA 5106}

%
% Title Page
%

\title{
    \vspace{2in}
    \textmd{\textbf{\hmwkClass:\ \hmwkTitle}}\\
    \normalsize\vspace{0.1in}
}

\author{Yousef Alaa Awad}

\begin{document}

\maketitle
\pagebreak

\section{Tomasulo's Out-of-Order Execution}

\textbf{Problem Setup:}
\begin{itemize}
	\item \textbf{A:} \texttt{load F1, 0 (R0)} (ROB1)
	\item \textbf{B:} \texttt{load F2, 0 (R1)} (ROB2)
	\item \textbf{C:} \texttt{subd F3, F1, F2} (ROB3)
	\item \textbf{D:} \texttt{addi F4, F2, \#5} (Not Dispatched)
	\item \textbf{E:} \texttt{multd F5, F3, F2} (Not Dispatched)
\end{itemize}

\textbf{Conditions:}
\begin{enumerate}
	\item Prior instructions have retired.
	\item Instructions A, B, and C have dispatched.
	\item Instruction A has completed (MEM[2500] = 198) but not retired.
\end{enumerate}

\subsection{(a) Register File}
The following table shows the state of the Register File.
Registers that have been renamed but not retired are marked as not being in the RF, and their corresponding ROB tags are provided.
When a destination register has been renamed but the instruction has not retired yet, the Value column keeps the old committed RF value.

\begin{table}[H]
	\centering
	\renewcommand{\arraystretch}{1.3}
	\begin{tabular}{|l|c|c|c|}
		\hline
		\textbf{Index} & \textbf{In RF?} & \textbf{Tag} & \textbf{Value} \\
		\hline
		R0             & Y               & B            & 2500           \\
		\hline
		R1             & Y               & B            & 2500           \\
		\hline
		F0             & Y               & B            & 5              \\
		\hline
		F1             & N               & ROB1         & 10             \\
		\hline
		F2             & N               & ROB2         & 15             \\
		\hline
		F3             & N               & ROB3         & 20             \\
		\hline
		F4             & Y               & B            & 25             \\
		\hline
		F5             & Y               & B            & 30             \\
		\hline
	\end{tabular}
\end{table}

\subsection{(b) Reservation Stations}
The state of the Reservation Stations for instructions B (ROB2) and C (ROB3) is shown below. Note that for instruction C, the first operand (F1) is ready because instruction A (ROB1) has completed and broadcast its result.

\begin{table}[H]
	\centering
	\renewcommand{\arraystretch}{1.3}
	\begin{tabular}{|l|c|c|c|c|}
		\hline
		\textbf{Tag} & \textbf{S1 Ready?} & \textbf{S1 Tag/Value} & \textbf{S2 Ready?} & \textbf{S2 Tag/Value} \\
		\hline
		ROB2         & y                  & 2500                  & y                  & 0                     \\
		\hline
		ROB3         & y                  & 198                   & n                  & ROB2                  \\
		\hline
	\end{tabular}
\end{table}

\end{document}
